% Part: second-order-logic
% Chapter: syntax-and-semantics
% Section: semantic-notions

\documentclass[../../../include/open-logic-section]{subfiles}

\begin{document}

\olfileid{sol}{syn}{sem}
\olsection{चिन्हार्थविषयक संकल्पना}

\begin{explain}
\emph{वैधता},
\emph{तार्किक निष्पन्नता}
आणि
\emph{पूर्ततायोग्यता}
या मध्यवर्ती
तार्किक संकल्पना
द्वितीय-क्रम
तर्कशास्त्रासाठी
प्रथम-क्रम
तर्कशास्त्राप्रमाणेच
परिभाषित होतात;
फरक इतकाच की
आधारभूत पूर्ति-संबंध
आता द्वितीय-क्रम
!!{formula}s
साठीचा असतो.
अर्थात, द्वितीय-क्रम
!!{sentence}
हे असे
!!a{formula}
असते ज्यातील
विधेय-चल आणि
फलन-चल यांसह
सर्व चले
बद्ध असतात.
\end{explain}

\begin{defn}[वैधता]
वाक्य $!A$
\emph{वैध},
$\Entails !A$,
असते जर आणि
फक्त जर प्रत्येक
!!{structure}~$\Struct M$
साठी $\Sat{M}{!A}$.
\end{defn}

\begin{defn}[तार्किक निष्पन्नता]
वाक्यांच्या
संच~$\Gamma$
मधून वाक्य~$!A$
\emph{तार्किकरीत्या
निष्पन्न होते},
$\Gamma
\Entails !A$,
जर आणि फक्त जर
$\Sat{M}{\Gamma}$
असलेल्या प्रत्येक
!!{structure}~$\Struct M$
साठी
$\Sat{M}{!A}$.
\end{defn}

\begin{defn}[पूर्ततायोग्यता]
एखाद्या
!!{structure}~$\Struct M$
साठी
$\Sat{M}{\Gamma}$
असेल, तर वाक्यांचा
संच~$\Gamma$
\emph{पूर्ततायोग्य}
असतो.
$\Gamma$
पूर्ततायोग्य नसेल,
तर त्याला
\emph{अपूर्ततायोग्य}
म्हणतात.
\end{defn}

\end{document}
